%%%% APÊNDICE B

\chapter{Comandos de Construção, Teste e Execução}\label{cap:apendiceb}

Este apêndice reúne comandos úteis para reproduzir a construção e validação do protótipo.

\section{Validação local}\label{sec:apendiceValidacaoLocal}

A \autoref{codigo:comandosTeste} reúne os alvos de validação local usados no trabalho.

\begin{sourcecode}[htb]
\caption{\label{codigo:comandosTeste}Comandos de teste local}
\begin{lstlisting}[frame=single]
make docker-image
make cpp-test
make cpp-smoke
make vhdl-test-all
\end{lstlisting}
\fonte{Elaborado pelos autores}
\end{sourcecode}

O alvo \codigo{cpp-test} valida o invólucro MMIO com arquivo temporário. O alvo \codigo{cpp-smoke} executa o gerador e o receptor em modo software. O alvo \codigo{vhdl-test-all} executa os bancos de teste da ponte, do livro, da estratégia, do motor completo e do invólucro Avalon-MM.

\section{Geração da estratégia}\label{sec:apendiceGeracaoEstrategia}

A \autoref{codigo:geracaoHdl} mostra os comandos de preparação do ambiente MATLAB e geração da estratégia.

\begin{sourcecode}[htb]
\caption{\label{codigo:geracaoHdl}Geração do bloco de estratégia por HDL Coder}
\begin{lstlisting}[frame=single]
make matlab-docker-image
make matlab-hdl-generate
\end{lstlisting}
\fonte{Elaborado pelos autores}
\end{sourcecode}

O primeiro comando prepara a imagem esperada para execução do MATLAB em ambiente controlado. O segundo executa a rotina \codigo{matlab/hdl_generator.m}, que chama a geração VHDL da função \codigo{strategy.m}. Após essa etapa, recomenda-se executar os bancos de teste da estratégia e do motor completo, pois uma mudança no bloco gerado pode alterar latência, portas ou comportamento.

\section{Geração HDL e compilação para DE10-Nano}\label{sec:apendiceDe10}

A \autoref{codigo:comandosDe10} apresenta a sequência de preparação, compilação, programação e implantação na placa.

\begin{sourcecode}[htb]
\caption{\label{codigo:comandosDe10}Fluxo previsto para DE10-Nano}
\begin{lstlisting}[frame=single]
make de10-sysroot
make matlab-login
make build
make quartus-program
make de10-enable-bridges
make deploy
\end{lstlisting}
\fonte{Elaborado pelos autores}
\end{sourcecode}

Após a programação do FPGA, o gerador e o receptor podem ser executados na placa com os comandos da \autoref{codigo:execucaoDe10}:

\begin{sourcecode}[htb]
\caption{\label{codigo:execucaoDe10}Execução na placa}
\begin{lstlisting}[frame=single]
ssh root@192.168.7.1 'cd /home/root && ./fast_data_feed'
ssh root@192.168.7.1 \
  'cd /home/root && HFT_FPGA_MMIO_BASE=0xFF200000 ./fast_receiver'
\end{lstlisting}
\fonte{Elaborado pelos autores}
\end{sourcecode}

O endereço \codigo{0xFF200000} corresponde à base típica da ponte leve HPS--FPGA na DE10-Nano quando o componente está no offset zero do Platform Designer. Se o componente for posicionado em outro offset, a base deve ser ajustada.

\section{Variáveis de ambiente}\label{sec:apendiceVariaveis}

A \autoref{tab:apendiceVariaveis} relaciona as variáveis que habilitam e configuram o acesso físico. Na ausência da variável de base, o receptor não tenta mapear o FPGA.

\begin{table}[htb]
\caption{Variáveis de ambiente do receptor}
\label{tab:apendiceVariaveis}
\begin{tabularx}{\textwidth}{lX}
\toprule
\textbf{Variável} & \textbf{Função} \\ \midrule
\codigo{HFT_FPGA_MMIO_BASE} & Base física do periférico. Quando ausente, o receptor opera somente em software. \\
\codigo{HFT_FPGA_MMIO_SPAN} & Tamanho da região mapeada. Quando ausente, usa valor padrão suficiente para a geometria esperada. \\
\codigo{HFT_FPGA_MMIO_DEV} & Dispositivo usado para mapeamento MMIO. Por padrão, \codigo{/dev/mem}. \\ \bottomrule
\end{tabularx}
\fonte{Elaborado pelos autores conforme \codigo{fast_receiver.cpp}}
\end{table}

A presença de \codigo{HFT_FPGA_MMIO_BASE} muda o modo de operação do receptor. Por isso, testes de software devem ser executados sem essa variável, enquanto testes em placa devem defini-la explicitamente. Em ambientes compartilhados, recomenda-se imprimir a base efetiva antes de enviar eventos para evitar confusão entre offset de Platform Designer e base da ponte leve HPS--FPGA.

\section{Sequência recomendada de depuração}\label{sec:apendiceDebug}

Uma sequência prática de depuração é:

\begin{enumerate}
  \item executar \codigo{make cpp-test} para validar a semântica de filas no invólucro;
  \item executar \codigo{make cpp-smoke} para validar FAST, TCP e decodificação sem FPGA;
  \item executar \codigo{make vhdl-test-all} para validar os blocos de hardware em simulação;
  \item gerar e programar o arquivo de configuração da FPGA;
  \item habilitar pontes HPS--FPGA;
  \item executar o receptor com \codigo{HFT_FPGA_MMIO_BASE} e verificar leitura de \codigo{MAGIC};
  \item iniciar o gerador e observar respostas \codigo{[FPGA->ARM]}.
\end{enumerate}

Se a leitura de \codigo{MAGIC} falhar, o problema provavelmente está em endereço, arquivo de configuração ou ponte HPS--FPGA. Se \codigo{MAGIC} estiver correto, mas \codigo{Send} falhar por fila cheia, deve-se investigar o consumo da fila TX, a reinicialização e os sinais internos de validade e prontidão. Se respostas não aparecerem, mas TX for consumida, o foco passa a ser o livro de ofertas, a estratégia e a fila RX.
